\subsection{Putting it all together}\label{sec:final_algo}
We conclude the section by putting all the results related to the procedures involved in Algorithm~\ref{alg:final_algorithm} together, in order to derive the guarantees of the algorithm.


% Finally, we show how to employ the \texttt{Action-Oracle} and \texttt{Try-Partition} algorithms to learn an approximately optimal contract with $\mathcal{O}(m^n)$ samples. Firstly, we observe that the number of iterations required by the while loop at Lines in the \texttt{Discover-and-Partition} algorithm is bounded by $2n$. This is because at each iteration, we either increase the size of the set of meta-actions $\dreg$ by one or decrease its size by merging one or more elements. Formally:

First, by Lemma~\ref{lem:finaliterations} the number of calls to the \texttt{Try-Cover} procedure is at most $2n$.
%
%Intuitively, this is because the \texttt{Try-Cover} procedure is re-started only when the previous call ended prematurely with $\varnothing$ as return value.
%
%This occurs only if some call to \texttt{Action-Oracle} returns the special value $\bot$, which may happen at most $2n $ thanks to Lemma~\ref{lem:finaliterations}.
%
%
%if the \texttt{Try-Cover} procedure terminates prematurely with $\varnothing$ and needs to be re-started, then the size of the set $\mathcal{D}$ of meta-actions has either increased by one or decreased by at least one (when two or more meta-actions are merged together).
%%
%Both cases happen at most $n$ times.
%%
%Formally:
%%
%\begin{comment}
%\begin{restatable}{lemma}{finaliterations}\label{lem:finaliterations}
%	Under the event $\mathcal{E}_\epsilon$, Algorithm~\ref{alg:final_algorithm} calls \textnormal{\texttt{Try-Partition}} at most $2n$ times.
%	%
%	%The number of iterations required \textnormal{\texttt{Discover-and-Partition}} is of at most $2n$.
%\end{restatable}
%\end{comment}
%
%
Moreover, by Lemma~\ref{lem:part_compl} and by definition of $q$ and $\eta$, the number of rounds required by each call to \texttt{Try-Cover} is at most $\widetilde{\mathcal{O}}( m^n \cdot \mathcal{I} \cdot \nicefrac{1}{\rho^4} \log (\nicefrac{1}{\delta})
)$ under the event $\mathcal{E}_\epsilon$, where $\mathcal{I}$ is a term that depends polynomially in $m$, $n$, and $B$.
%
Finally, by Lemma~\ref{lem:solve_lpl} the contract returned by Algorithm~\ref{alg:final_algorithm}---the result of a call to the \texttt{Find-Contract} procedure---has expected utility at most $\rho$ less than the best contact in $[0,B]^m$,
%
%$\mathcal{O}(\sqrt{\epsilon})$-optimal under the clean event.
%
while the probability of the clean event $\mathcal{E}_\epsilon$ can be bounded below by means of a concentration argument.
%
All the observations above allow us to state the following main result.
%
%Next, by means of Lemma~\ref{lem:solve_lpl}, we ensure that the solution returned by the LPs is $\mathcal{O}(\sqrt{\epsilon})$-optimal. Moreover, at each iteration, the number of iterations required by the \texttt{Try-Partition} algorithm is at most $\mathcal{O}(m^n)$ by Lemma~\ref{}. Finally, exploiting this information and noting that the probability of the clean event $\mathcal{E}_\epsilon$ can be bounded using a union bound and a concentration argument, we can show the following theorem.
%
\begin{restatable}{theorem}{finaltheorem}\label{thm:finalthm}
	Given $\rho\in (0,1)$, $\delta \in (0,1)$, and $B \geq 1$ as inputs, with probability at least $1-\delta$ the \textnormal{\texttt{Discover-and-Cover}} algorithm (Algorithm~\ref{alg:final_algorithm}) is guaranteed to return a contract $p \in [0,B]^m$ such that $u(p) \geq \max_{p' \in [0,B]^m} u(p') - \rho$ in at most $\widetilde{\mathcal{O}}( m^n \cdot \mathcal{I} \cdot \nicefrac{1}{\rho^4} \log(\nicefrac{1}{\delta})	)$ rounds, where $\mathcal{I}$ is a term that depends polynomially in $m$, $n$, and $B$.
	%
	% Given $\alpha,\epsilon >0$  with $O(m^m)$ samples \textnormal{\texttt{Discover-and-Partition}} returns a ..-optimal solution with probability of at least 1-...
\end{restatable}
%Notice that the bound on the number of rounds provided by Theorem~\ref{thm:finalthm} is polynomial in the instance size (including the number of outcomes $m$) when the number of agent's actions $n$ is constant.
Notice that the number of rounds required by Algorithm~\ref{alg:final_algorithm}is polynomial in the instance size (including the number of outcomes $m$) when the number of agent's actions $n$ is constant.

%\alb{Ho usato la stessa complexity per la try partition singola, va controllato.}\bac{cioe?}

%\bac{la dipendenza nel bound da $n$ è omessa dato che è considerata costante ... }
%\alb{$\epsilon$-optimal non va bene, bisogna scrivere la garanzia di LOBC.}
%
%\alb{Forse per come diamo le garanzie non ha senso definire LOBC come problema, decidiamo.}

%We conclude the section with some remarks on the $\texttt{Discover-and-Partition}$ algorithm. First, the $\texttt{Discover-and-Partition}$ can be seen as a generalization of the algorithm designed by~\cite{peng2019optimal} for learnng optimal strategies in repeated Stackelberg games. The challenging task is that in our setting we do not observe the action played by the agent and thus we have the necessity of designing an suitble defined algorithm,namely \texttt{Action-Oracle}. Furtheremore, we propoe an algorithm exploitig the  

%\alb{Questo va in related works:
%We conclude the section outlining that our algorithm can be seen as a generalization of the one developed by \cite{peng2019optimal} for learning optimal strategies in repeated Stackelberg games. However, our setting poses an additional challenge as we do not observe the action played by the agent, and thus, we design a suitably defined algorithm, namely the \texttt{Action-Oracle} to overcame this issue.}
\section{Connection with online learning in principal-agent problems}\label{sec:regret}
In this section, we show that our \texttt{Discover-and-Cover} algorithm can be exploited to derive a no-regret algorithm for the related online learning setting in which the principal aims at maximizing their cumulative utility.
%
In such a setting, the principal and the agent interact repeatedly over a given number of rounds $T$, as described in Section~\ref{sec:learning_problem}.

\begin{wrapfigure}[10]{R}{0.46\textwidth}
	\vspace{-1.3cm}
	\begin{minipage}{0.44\textwidth}
		\begin{algorithm}[H]
			\caption{No-regret algorithm}
			\label{alg:regret}
			\small
			\begin{algorithmic}[1]
				\Require $\delta \in (0,1)$, $B \geq 1$  
				\State Set $\rho$ as in proof of Theorem~\ref{thm:regret}
				\For {$t=1,\ldots,T$}
				\If {Algorithm~\ref{alg:final_algorithm} \emph{not} terminated yet}
				\State $p^t \gets p \in \mathcal{P}$ prescribed by Alg.~\ref{alg:final_algorithm}
				\Else
				\State $p^t \gets p\in [0,B]^m$ returned by Alg.~\ref{alg:final_algorithm}
				\EndIf
				\EndFor
				%		
				%		\State \bac{non so come dirlo meglio questo while} 
				%		\While{Algorithm~\ref{alg:final_algorithm} is not ended} 
				%		\State Commit to contract $p$ prescribed by Algorithm~\ref{alg:final_algorithm} \Comment{Exploration Phase}
				%		\EndWhile
				%		\State Commit to $p^* \in \preg$ prescribed by Algorithm~\ref{alg:final_algorithm} \Comment{Exploitation Phase}
			\end{algorithmic}
		\end{algorithm}
	\end{minipage}
\end{wrapfigure}

At each $t = 1, \ldots, T$, the principal commits to a contract $p^t \in \mathbb{R}_+^m$, the agent plays a best response $a^\star(p^t)$, and the principal observes the realized outcome $\omega^t \sim F_{a^\star(p^t)}$ with reward $r_{\omega^t}$.
%
% As customary in the literature on online learning, at each round $t$ the principal selects a contract to commit by means of a (possibly randomized) policy $\pi^t : \mathcal{W}^{t-1} \to \mathbb{R}_+^m$, which maps any history of the interaction $\mathcal{W}^{t-1} :=\{p^{1}, \omega^{1}, \dots, p^{t-1}, \omega^{t-1}\}$ up to round $t-1$ to a probability distribution over contracts.
%
Then, the performance in terms of cumulative expected utility by employing the contracts $\{p^t\}_{t=1}^T$ is measured by the \emph{cumulative (Stackelberg) regret}
%
%
% In this section, we formally introduce the concept of \emph{cumulative regret}. We first define the set of policies $\pi_t : \mathcal{W}^{(t-1)} \to \preg$ as (possibly randomized) maps from the history of the interaction $\mathcal{W}^{(t-1)}=\{p^{1}, \omega^{1}, \dots, p^{t-1}, \omega^{t-1}\}$ to the set of contracts $\preg$, specifying at each time $t \in [T]$ a probability distribution over the set of contracts that the principal may commit to. We define the \emph{expected regret} suffered by a family of policies $\pi \coloneqq \{\pi_t\}_{t=1}^T$ as follows:
%
%\begin{small}
%	\begin{equation*}
		%\hspace{-60mm}
%		$$
%		R^T := T \cdot\max_{p \in [0,B]^m} u(p) - \mathbb{E} \left[\sum_{t=1}^T    u(p^t) \right] ,$$
		$$
		R^T := T \cdot\max_{p \in [0,B]^m} u(p) - \sum_{t=1}^T    u(p^t)  ,$$
%	\end{equation*}
%\end{small}
%
%where the expectation is over the randomness of the environment.
%which is determined by the stochasticity of the realized outcome in our setting. 
%
%such as the outcome sampled from the distributions associated with the agent's actions.
As shown in Section~\ref{sec:discover_and_partition}, Algorithm~\ref{alg:final_algorithm} learns an approximately-optimal bounded contract with high probability by using the number of rounds prescribed by Theorem~\ref{thm:finalthm}.
%
% $\mathcal{O}(m^n)$ \bac{$\mathcal{O}(\mathcal{I} m^n)$} rounds.
%
Thus, by exploiting Algorithm~\ref{alg:final_algorithm}, it is possible to design an \emph{explore-then-commit} algorithm ensuring sublinear regret $R^T$ with high probability; see Algorithm~\ref{alg:regret}.
%
%by employing $\mathcal{O}(m^n)$ rounds, it is possible to learn an approximately optimal contract with a high probability. Consequently, we can design an "explore-then-commit" algorithm ensuring sublinear regret with a high probability. Please refer to Algorithm~\ref{alg:regret} for the detailed steps. 
%
%By Theorem~\ref{thm:finalthm}, we can prove the following:
%Then, the following holds:
%
\begin{restatable}{theorem}{finalregret}\label{thm:regret}
Given $\alpha \in (0,1)$, Algorithm~\ref{alg:regret} achieves $R^T  \le \widetilde{\mathcal{O}} \left(  m^n \cdot \mathcal{I} \cdot  \log (\nicefrac{1}{\delta}) \cdot 
T^{4/5} \right)$ with probability at least $1-\delta$, where $\mathcal{I}$ is a term that depends polynomially on $m$, $n$, and $B$.
%
%the \emph{cumulative regret} suffered by Algorithm~\ref{alg:regret} is
%\begin{equation*}
%R^T  \le \widetilde{\mathcal{O}} \left( B m^{n}\log\left({2mT}/{\alpha}\right) T^{4/5} \right),
%\end{equation*}
%with probability at least $1-\alpha$.
\end{restatable}
%
Notice that, even for a small number of outcomes $m$ (\emph{i.e.}, any $m \geq 3$), our algorithm achieves better regret guarantees than those obtained by~\citet{zhu2022sample} in terms of the dependency on the number of rounds $T$.
%
% , when the number of agent's actions $n$ is a constant.
%
Specifically, \citet{zhu2022sample} provide a $\widetilde{\mathcal{O}}(T^{1-1/(2m+1)})$ regret bound, which exhibits a very unpleasant dependency on the number of outcomes $m$ at the exponent of $T$.
%
Conversely, our algorithm always achieves a $\widetilde{\mathcal{O}}(T^{4/5})$ dependency on $T$, when the number of agent's actions $n$ is small.
%
This solves a problem left open in the very recent paper by~\citet{zhu2022sample}.
%
% when the contract space is $[0,1]^m$, while our algorithm achieves regret $\widetilde{\mathcal{O}}(T^{4/5})$ when $n$ is 
%
% Then, our algorithm outperforms theirs in terms of regret bound when the number of actions is constant.
%
%\bac{droppare sta parte}
%Additionally, our algorithm scales with the parameter $B$, which represents the bounding space of contracts. This scalability results in a new dimension of flexibility in contract design, allowing for larger contract spaces while still achieving competitive regret guarantees.

%As a final remark, it is worth noting that our algorithm, even with $m=1$, achieves comparable regret guarantees of the ones probosed by~\cite{zhu2022sample}. In their work, they provide a regret bound of $\mathcal{O}(T^{(2m+2)/(2m+3)})$, matching our . Moreover, our algorithm outperforms theirs in terms of regret bound when the number of actions is fixed. Additionally, our algorithm scales with the parameter $B$, which represents the bounding space of contracts. This scalability results in a new dimension of flexibility in contract design, allowing for larger contract spaces while still achieving competitive regret guarantees.

%\alb{La parte sopra la discuterei insieme.}
%
%\alb{Controllare la baseline dello Stackelberg regret.}



\begin{comment}
\subsubsection*{Event}

$\mathcal{E}_\epsilon$ is the event in which all the calls to \texttt{ACTION-ORACLE} return an $\tilde F$ such that $ \|\widetilde{F} - F_{a^\star(p)} \|_{\infty} \le \epsilon $, where $p \in \mathcal{P}$ is the contract given as input to the procedure.
%
Intuitively, \texttt{ACTION-ORACLE} works properly.

\subsubsection*{\texttt{TRY-PARTITION}}

Let $\mathcal{P}_d \subseteq \mathcal{P}$ be the set of contracts $p \in \mathcal{P}$ given as input to \texttt{ACTION-ORACLE} all the times it returned an $\tilde F$ such that $\tilde F \in \mathcal{F}_d$ (at the end of the algorithm).


Let $I(d) \coloneqq \bigcup_{p \in \mathcal{P}_d} a^\star(p)$.

\begin{lemma}
	Under $\mathcal{E}_\epsilon$, if $\tilde F = \bot$, for every discovered action $d \in \mathcal{D}$, contract $p \in \mathcal{L}_{d}$, and action $a_i \in I(d)$, it holds:
	\begin{equation*}
		\sum_{\omega \in \Omega } p_{\omega}F_{a, \omega} - c_{a} \ge
		\sum_{\omega \in \Omega } p_{\omega}F_{a', \omega} - c_{a'} -  \epsilon',
	\end{equation*}
	for every action $a' \in \mathcal{A}$, where $\epsilon' = 14 \epsilon m n$.
\end{lemma}

\begin{lemma}
	Under $\mathcal{E}_\epsilon$, if $\tilde F = \bot$, it holds:
	\[
	\bigcup_{d \in \mathcal{D}} \mathcal{L}_d = \widetilde{\mathcal{P}}.
	\]
\end{lemma}


\begin{lemma}
	Under $\mathcal{E}_\epsilon$, if $\tilde F \neq \bot$, it holds that $\left\{ d \in \mathcal{D} \mid \exists F \in \mathcal{F}_d :  \|{F} - \widetilde{F} \|_{\infty} \le 2\epsilon \right\} \neq 1$.
\end{lemma}


\subsubsection*{\texttt{UPDATE-ACTIONS}}

\begin{lemma}
	Algorithm~\ref{alg:action_oracle} ensures that for each discovered action $d \in \mathcal{D}$ and for each actions $a_i, a_j \in I(d)$ we have $\| F_{a_i} -   F_{a_j} \|_{\infty} \le 5 \epsilon n$. Furthermore for each empirical distribution $\widetilde{F} \in \mathcal{F}_{d}$ we have: $\|\widetilde{F}- F_{a_i} \|_{\infty} \le 5 \epsilon n $ for each  $a_i \in I(d)$.
\end{lemma}

\begin{lemma}
	$I(d) \cap I(d') = \varnothing$ for every $d,d' \in \mathcal{D}$.
\end{lemma}

\subsubsection*{\texttt{FIND-CONTRACT}}

\begin{lemma}
	
\end{lemma}

\hrulefill
\end{comment}
%\clearpage
%\hrulefill
